\documentclass[11pt]{article}
\usepackage[margin=1in]{geometry}
\usepackage{amsmath,amssymb,amsthm}
\usepackage[hidelinks,pdfauthor={Bowen Zuo},pdftitle={Interlacing of excedance polynomials with a minimum cycle length}]{hyperref}
\newtheorem{theorem}{Theorem}
\newtheorem{lemma}[theorem]{Lemma}
\newcommand{\exc}{\operatorname{exc}}
\newcommand{\ii}{\mathrm{i}}
\allowdisplaybreaks[2]
\title{Interlacing of excedance polynomials with a minimum cycle length}
\author{Bowen Zuo\\[2pt]\small East China University of Science and Technology}
\date{November 2026}
\begin{document}
\maketitle
\begin{abstract}
For every fixed cutoff \(k\ge3\), the excedance polynomials of derangements with independent weights in \([0,1]\) on cycles of lengths two through \(k-1\) and unit weights on longer cycles eventually interlace strictly after removing their common simple zero. The size threshold is uniform in the weights. We prove the unweighted result at cutoff four for every nondegenerate pair and give an independent proof at cutoff three. At both cutoffs the corresponding weighted families interlace from upper size 512. The proofs use comparisons at genuine predecessor roots and exact finite certificates. All four theorems are formalized in Lean~4.
\end{abstract}

\section{The statements}
For a permutation \(\sigma\) on \([j]\), put \(\exc(\sigma)=\#\{i:i<\sigma(i)\}\), and let \(m_2(\sigma)\) be its number of cycles of length two. Set
\[
 F_j(t,u)=\sum_{\substack{\sigma\in\mathfrak S_j\\ \sigma(i)\ne i\ \forall i}}
       t^{\exc(\sigma)}u^{m_2(\sigma)},\qquad
 P_j(t,u)=F_j(t,u)/t\quad(j\ge3).
\]
Taking \(u=0\) retains exactly the permutations with every cycle of length at least three. Alexandersson asks about real-rootedness and consecutive interlacing for every fixed minimum cycle length \(k\ge3\) \cite{alex}. The recurrence for these rows follows also by summing the finer enumeration of Rakotoson and Rakotondrajao \cite[Theorem~30]{rr}. Brenti proves real-rootedness on individual conjugacy classes and unimodality results for restrictions on cycle lengths \cite[Theorems~3.2--3.4]{brenti}.
For a cutoff \(k\), let \(E_{j,k}(t)\) denote the sum of \(t^{\exc(\sigma)}\) over permutations whose cycles all have length at least \(k\).
In contemporaneous work, ProgramComputer \cite{programcomputer} proves the unweighted \(k=3\) case at every size and gives absolute bounds \(n\ge Bk^2\) and \(n\ge10^{20}\) for eventual unweighted interlacing, where \(n\) is the lower size. The latter bound works simultaneously for all admissible cutoffs. Their two-cycle theorem covers the initial pairs through upper size \(3k-1\). Our results concern independent cycle weights, all sizes at cutoff four, and formal verification. We include the independent unweighted cutoff-three proof in the same Lean development.
\begin{theorem}\label{general}
For every fixed integer \(k\ge3\) there is an \(N(k)>k\) such that, for all \(n\ge N(k)\), the consecutive quotients \(E_{n-1,k}/t,E_{n,k}/t\) have respectively \(n-3,n-2\) simple negative roots. Between each two consecutive roots of the upper row lies exactly one root of the lower row. Each original row has one further simple root at zero, and the full rows weakly interlace.

The threshold is uniform over all choices of weights \(w_2,\ldots,w_{k-1}\in[0,1]\), when each derangement receives the product of its short-cycle weights, with unit weight for each longer cycle.
\end{theorem}
The proof of Theorem~\ref{general} is in Section~\ref{sec:general}. For minimum cycle length three, the finite range below the analytic threshold can be completed.
\begin{theorem}\label{main}
For every \(n\ge4\), the roots of \(P_{n-1}(t,0)\) and \(P_n(t,0)\) are simple and negative. If the upper row has roots \(\beta_1<\cdots<\beta_{n-2}\) and the preceding row has roots \(\alpha_1<\cdots<\alpha_{n-3}\), then \(\beta_i<\alpha_i<\beta_{i+1}\) for all \(1\le i\le n-3\). Both full polynomials have one further simple root at zero and weakly interlace.
\end{theorem}
\begin{theorem}\label{family}
For \(n\ge512\), all the conclusions of Theorem~\ref{main} hold with any common mark \(u\in[0,1]\) in the two rows.
\end{theorem}
\begin{theorem}\label{four}
For minimum cycle length four the consecutive quotients \(E_{n-1,4}/t,E_{n,4}/t\) strictly interlace with simple negative roots at every upper index \(n\ge5\). The full rows weakly interlace. For \(n\ge512\) the same conclusions hold when derangements receive independent weights \(w_2,w_3\in[0,1]\) for their cycles of lengths two and three, respectively, and unit weight for every longer cycle.
\end{theorem}
We prove Theorem~\ref{family} by a uniform sign test at actual roots. Section~\ref{finite} completes Theorem~\ref{main} using finite integer inequalities. The coefficient of \(t\) in \(F_j\) is one: every cycle contributes at least one excedance and there is exactly one \(j\)-cycle having just one excedance. For example, the latter assertion follows from the single-cycle polynomial \(tA_{j-1}(t)\), with \(A_m\) the ordinary Eulerian polynomial. Inverting a fixed-point-free permutation sends each excedance to a deficiency and preserves its cycle type. Thus \(P_j\) is reciprocal of degree \(j-2\), monic, has constant term one, and is positive on \([0,\infty)\).

\section{Transport of a root sign}\label{transport}
Let \(P,Q\) be monic of degrees \(d,d+1\). Assume their coefficients vary continuously on an interval and both polynomials are positive on \([0,\infty)\). Suppose at every real root \(r\) of \(P\) we have \(Q(r)P'(r)<0\). If \(P\) splits at one parameter, it splits throughout: splitting is closed, since the ordered real roots of a convergent sequence of monic polynomials remain bounded and a subsequential limit gives a full factorization. It is open here because every real root is simple and a disjoint sign-change interval around each persists under sufficiently small coefficient changes. At any parameter the signs of \(Q\) alternate at the \(d\) ordered roots. The signs at \(-\infty\) and at zero give the two exterior roots, hence all \(d+1\) roots, with multiplicity one.

At \(u=1\) the ordinary derangement rows supply the splitting anchor. Define the auxiliary initial row \(P_2(t,1)=F_2(t,1)/t=1\). Then \(P_3(t,1)=1+t\), and for \(n\ge4\) their insertion recurrence is
\[
 P_n=(1+(n-2)t)P_{n-1}+t(1-t)P'_{n-1}+(n-1)tP_{n-2}.
\]
If \(P_{n-2}\) strictly interlaces \(P_{n-1}\), its sign at any root of \(P_{n-1}\) equals the sign of the derivative there. The two surviving terms thus both have sign opposite to that derivative, giving strict interlacing by induction.

For \(t=-e^{-x}\) write \(H_j=F_j(t,u)/(1-t)^j\), and let \(p=n-1\). Primes on \(H_p\) and \(P_p\) denote differentiation with respect to \(x\) and \(t\), respectively. At a predecessor root, \(-H'_p=t^2P'_p(t,u)/(1-t)^p\). The sign
\begin{equation}\label{goal}
 H_n(-H'_p)>0\qquad(H_p=0,\quad x\ge0)
\end{equation}
therefore implies the required polynomial sign on \([-1,0)\). If \(t<-1\), reciprocal symmetry gives \(P'_p(t)=-t^{p-4}P'_p(1/t)\) at the root and \(P_n(t)=t^{n-2}P_n(1/t)\). Their product changes by a positive factor. Thus \eqref{goal} and the anchor imply Theorem~\ref{family}.

\section{The normalized marked row}
The exponential generating function for all permutations by excedances is \((1-t)/(\exp((t-1)v)-t)\). By the labelled cycle formula, deleting fixed points multiplies it by \(e^{-v}\); marking two-cycles multiplies it by \(\exp((u-1)tv^2/2)\). Substitute \(v=z/(1-t)\), and set
\[
 L(z)=(1+e^{z-x})^{-1},\quad B(z)=e^{az+bz^2},\qquad
 a=\frac{-t}{1-t},\quad b=\frac{(u-1)t}{2(1-t)^2}.
\]
Then \(G(z)=\sum H_j z^j/j!=(1-t)B(z)L(z)\).
The bounds \(0\le a\le\min(1/2,e^{-x})\) and \(0\le b\le\min(1/8,e^{-x}/2)\) hold for \(x\ge0\). Since \((\partial_x+\partial_z)L=0\), direct differentiation gives
\begin{align}
 H_n+H'_p&=p\alpha H_{n-2}+p(p-1)\beta H_{n-3},\label{insert}\\
 \alpha&=ut/(1-t)^2,\qquad \beta=(1-u)t(1+t)/(2(1-t)^3).\nonumber
\end{align}
In particular \(|\alpha|,|\beta|\le e^{-x}\). Set
\[
 \zeta=x+\ii\pi=\rho e^{\ii\theta},\quad c=x/\rho,\quad s=\pi/\rho,\qquad
 Y_j=\frac{\rho^{j+1}H_j}{(1-t)j!},\qquad V_j=\rho^{j+1}[z^j]L(z).
\]
Write \(Z=\rho^{n+1}(H_n+H'_p)/((1-t)n!)\). Equation~\eqref{insert} reads
\[
 Z=(\alpha\rho^2Y_{n-2}+\beta\rho^3Y_{n-3})/n .
\]
The inequality \(|Z|<|Y_n|\) implies \eqref{goal} at any zero of \(H_p\).
\section{The three sign estimates}
Throughout this section \(n\ge512\) and \(0\le u\le1\).

\subsection{The first conjugate poles}
Assume \(0\le x\le16\), and put \(R=(x^2+4\pi^2)^{1/2}\). Then \(\rho<17,\ R<18,\ \rho/R<24/25\). The poles of \(L\) are \(x+\ii\ell\pi\), for odd signed integers \(\ell\), all with residue \(-1\). The contour \(|z|=R\) encloses precisely the two nearest ones. Its distance to each pole exceeds \(2/5\): the two smallest radial gaps are \(3\pi^2/(R+\rho)\) and \(5\pi^2/(R+\sqrt{x^2+9\pi^2})\). On the contour \(|1+e^{z-x}|>1/10\). Indeed, if \(v=\Re z-x\) has \(|v|\ge1\), use \(|1-e^v|\). Otherwise let \(w\in[-\pi,\pi]\) be the displacement of \(\Im z\) from the nearest odd multiple of \(\pi\). Then
\[
 |1+e^{z-x}|^2=(1-e^v)^2+4e^v\sin^2(w/2)
 \ \ge\ (v^2+w^2)/10.
\]
By \(aR\le xe^{-x}+\pi\) and \(bR^2\le x^2e^{-x}/2+\pi^2/2\), one has \(|B(z)|\le e^{17}\). At \(\zeta\), \(D=|B(\zeta)|\ge e^{-\pi^2/8}>1/4\). Moving the coefficient integral gives
\begin{equation}\label{poles}
 \left|Y_j-2\Re\{B(\zeta)e^{-\ii(j+1)\theta}\}\right|
 \le170\,e^{aR+bR^2}(\rho/R)^j <10^{-6}\qquad(j\ge509).
\end{equation}
For the last inequality split at \(x=8\). For \(x\le8\) one has \(\rho/R\le9/10\) by squaring, and \(170e^{17}(9/10)^{509}<10^{-6}\). For \(8\le x\le16\) the exponent is at most \((18+18^2/2)e^{-8}<1\), and \(170e(24/25)^{509}<10^{-6}\).

At \(Y_p=0\) let \(\phi=\arg B(\zeta)-n\theta\). Then \(|\cos\phi|\le2\cdot10^{-6}\) and \(|\sin\phi|>.99\). Since \(s\ge3/17\), the phase \(\phi-\theta\) in the next row and \eqref{poles} give \(|Y_n|>3Ds/2\). Both shorter rows have absolute value below \(3D\). For \(h=2,3\), \(\rho\le x+\pi\) and \(e^{-x}(x+\pi)^h\) decreases on the nonnegative half-line. Thus
\[
 |Z|\le 3D(\pi^2+\pi^3)/512 < (3D/2)(3/17).
\]

\subsection{The remaining central interval}
Assume \(16\le x\le3\sqrt{n-1}\). The partial fractions of the logarithmic derivative of the product for \(\cosh\) \cite{dlmf}, or its differentiated series, give
\[
 V_j=\sum_{\ell\ \mathrm{odd},\,\ell\in\mathbb Z}(c+\ii\ell s)^{-(j+1)}
 \quad (j\ge1),\qquad V_0=\rho/(1+e^{-x}).
\]
The differentiated sums converge absolutely and locally uniformly away from the poles. To control every order at once, write \(\ell=2i+1>0\). Each modulus of order at least two is at most \(\rho^2/(x^2+\pi^2\ell^2)\), which is at most
\[
 \rho^2\left(\frac1{x+2i}-\frac1{x+2i+2}\right),
\]
because \((x+\ell)^2-1\le2(x^2+\pi^2\ell^2)\). Telescoping and including conjugates gives \(|V_j|\le2\rho^2/x\le2(x+4)\) for \(j\ge1\); the constant term also satisfies the last bound. The nonnegative coefficients of \(B\) now give
\begin{align}
 |Y_j-V_j|&\le2(x+4)(e^{a\rho+b\rho^2}-1)
 \le4(x+4)^3e^{-x}<s/32.\label{bias}
\end{align}
Here \(a\rho+b\rho^2\le (x+4)^2e^{-x}<1/2\). For the last inequality use \(s\ge3/(x+4)\), the decrease of \(e^{-x}(x+4)^4\), and \(4\cdot20^4e^{-16}<3/32\).

For \(m\ge8,\ ms^2\ge1\) put \(U_m=V_{m-1}/2\). We claim
\begin{equation}\label{center}
 |U_m|\le1/8\quad\Longrightarrow\quad |U_{m+1}-cU_m|>s/2 .
\end{equation}
For positive odd \(\ell\), \(|c+\ii\ell s|^{-m}\le[1+(\ell^2-1)/8]^{-4}\), using \((1+h)^{m/8}\ge1+mh/8\). The sum of these bounds weighted by \(\ell\), over \(\ell\ge3\), is at most
\[
 3/16+4096(5^{-7}+5^{-6}/12)<1/3.
\]
Indeed the summand at 3 equals \(3/16\) in the bound and the remaining terms are at most \(4096\ell^{-7}\); compare the decreasing odd-lattice tail with an integral. The unweighted tail is below \(1/9\). Hence \(|\cos(m\theta)|\le1/8+1/9<1/4\). The \(\ell=1\) term of the centered difference has magnitude \(s|\sin(m\theta)|\). The remaining terms cost at most \(s/3\), since \(|(c+\ii\ell s)^{-1}-c|\le\ell s\) by squaring. This proves \eqref{center}.

In this interval \(ns^2\ge1\) since \(x^2\le9(n-1)\). At a zero of \(Y_p\), \eqref{bias} and \eqref{center} give \(|V_n-cV_p|>s\) and \(|Y_n|>15s/16\). Both shorter rows have absolute value below \(3(x+4)\), so
\[
 |Z|\le 6(x+4)^4e^{-x}/n<s/4.
\]
The last inequality uses the decrease of \(e^{-x}(x+4)^5\) and \(6\cdot 20^5e^{-16}/512<3/4\).

\subsection{The exterior}
Assume \(x\ge3\sqrt p\). Write \(r=e^{-x}\), \(w_h=h^p r^h\) for \(h\ge1\), and let \(m\) maximize these weights, with \(W=w_m\). If \(m>1\), comparison with \(m-1\) gives \(x\le p\log(m/(m-1))\le p/(m-1)\). Hence \(m\le2\sqrt p/5\) and \(m+1\le \sqrt p/3+2\le3\sqrt p/7\) for \(p\ge511\), also at \(m=1\). The product of the two neighbors divided by \(W^2\) equals \((1-m^{-2})^p\), with \(w_0=0\). The smaller neighbor is therefore below \(W/20\). Put \(\delta=1/200\). The ratio just past the right neighbor is at most \((1-(m+1)^{-2})^p\le e^{-49/9}<\delta\), by comparing successive ratios and using \(w_{m+1}\le W\). The left side has the same bound whenever that tail exists, and the ratios decrease outwards. Thus \(w_h\le W\delta^{|h-m|-1}\) at distances at least two. Summing yields
\[
 \sum_{h\ge1} w_h<4W,\qquad
 \sum_{h\ge1} h w_h\le m\sum w_h+\sum |h-m|w_h<6nW.
\]

Let \(T_j(h)=j![z^j]B(z)e^{hz}\). Expanding \(L\) geometrically near zero gives
\[
 H_j/(1-t)=\sum_{h\ge0}(-1)^h r^hT_j(h).
\]
For \(n-3\le j\le n\) and \(h\ge1\), expanding the tilt and bounding falling factorials by powers gives
\[
 1\le T_j(h)/h^j\le \exp(aj/h+bj^2/h^2)\le1+n^{-4}.
\]
Indeed \(x^2\ge9(n-1)\ge8n\) and \(32!\le8^{16}512^8\). Hence \(e^x\ge x^{32}/32!\ge n^8\) and \(r\le n^{-8}\). At \(h=0\) each monomial contributing to the coefficient uses at least \(\lceil j/2\rceil\) powers of \(r\); summing the inverse factorials of the linear and quadratic factors gives
\[
 0\le T_j(0)\le j!e^2 r^{\lceil j/2\rceil}\le n^{-4}r .
\]
For the last bound use \(\lceil j/2\rceil-1\ge j/3,\ x/3\ge2\log n\), and \(j!\le n^j\).
Put \(S_j=\sum_{h\ge1}(-1)^h r^hh^j\). As \(W\ge r\), the absolute error in replacing \(H_j/(1-t)\) by \(S_j\) is at most \(5n^{-4}W\) for \(n-3\le j\le p\) and at most \(7n^{-3}W\) for \(j=n\).

Let \(h_0\) be the label of the stronger immediate neighbor of \(m\), choosing the right one for \(m=1\). In \(S_n-h_0 S_p\) this neighbor disappears, while the mode contributes magnitude \(W\). Thus
\[
 |S_n-h_0 S_p|\ge W\left(1-2/20-2\sum_{d\ge2}(d+1)\delta^{d-1}\right)>4W/5 .
\]
At \(H_p=0\) the preceding error estimates and \(h_0\le n\) give \(|H_n|/(1-t)>3W/4\). Each of the two shorter rows has magnitude below \(5(1-t)W\); inserting these bounds in \eqref{insert} gives \(|H_n+H'_p|/(1-t)\le5(n+n^2)rW<W/4\). This proves \eqref{goal} also in the exterior and concludes the proof of Theorem~\ref{family}.

\section{Independent weights on the short cycles}\label{sec:general}
Fix \(k\ge3\), put \(d=k-1\), and take \(w_a\in[0,1]\) for \(2\le a<k\). Set \(w_1=0\) and \(w_a=1\) for \(a\ge k\). In this section \(F_j\) is the excedance polynomial of derangements with each cycle of length \(a\) contributing the additional factor \(w_a\). The same variables \(t=-e^{-x}\), \(p=n-1\), \(r=e^{-x}\), \(\zeta=\rho e^{\ii\theta}\), \(H_j,Y_j,V_j\) and \(Z\) as in the two-cycle case will be used with this \(F_j\).

\subsection{The change in successive cycle weights}
Define \(A_0(t)=1\) and, for \(a\ge1\),
\[
 M_a=\frac{tA_{a-1}(t)}{(1-t)^a},\qquad
 b_a=\frac{(w_a-1)M_a}{a!},\qquad
 B(z)=\exp\!\left(\sum_{a=1}^{d} b_a z^a\right).
\]
The labelled cycle formula gives \(\sum H_jz^j/j!=(1-t)B(z)L(z)\), with the same \(L(z)=(1+e^{z-x})^{-1}\). For \(a\ge2\) this follows by changing the single-cycle weight \(tA_{a-1}(t)\). The linear term includes deletion of fixed points: rewriting the unrestricted generating function in terms of \(L\) supplies the extra \(e^z\), and multiplication by \(e^{-z/(1-t)}\) leaves \(e^{-tz/(1-t)}\), exactly \(e^{b_1z}\). The ordinary Eulerian recurrence gives \(M'_a=-M_{a+1}\). Since \(w_{d+1}=1\), we set \(b_{d+1}=0\). It follows that
\begin{align}
 b'_a+(a+1)b_{a+1}&=(w_{a+1}-w_a)M_{a+1}/a! =: C_a,\\
 H_n+H'_p&=\sum_{a=1}^{d}(p)_a C_a H_{p-a},\qquad
 Z=\frac1n\sum_{a=1}^{d}C_a\rho^{a+1}Y_{p-a},\label{general-insertion}
\end{align}
where \((p)_a=p!/(p-a)!\) and \(p\ge d\). Indeed \((\partial_x+\partial_z)L=0\) and the constant term cancels as in \eqref{insert}; coefficient extraction gives both identities. In particular, when all short cycles have the same mark, only \(C_1\) and \(C_d\) can be nonzero. For \(x\ge0\), coefficient nonnegativity of the Eulerian polynomials gives
\[
 |A_{a-1}(-r)|\le A_{a-1}(1)=(a-1)!,\qquad
 |b_a|\le r/a,\quad |C_a|\le r.
\]
All estimates below are uniform over the weight cube for the fixed \(d\).

\subsection{A bounded interval}
For the moment fix \(X\ge16\). On \(0\le x\le X\), take \(R(x)=(x^2+4\pi^2)^{1/2}\). The only poles inside \(|z|=R(x)\) are \(\zeta,\bar\zeta\). By compactness, the integrand factor \(B(z)L(z)\) is bounded on these contours, uniformly over all weights; also \(\rho/R(x)\le q<1\). Moving the coefficient contour past both poles gives
\[
 Y_j=2\Re\!\left(B(\zeta)e^{-\ii(j+1)\theta}\right)+O_{d,X}(q^j).
\]
Here \(|B(\zeta)|\) has a strictly positive uniform lower bound, and \(s=\pi/\rho\) is bounded below. At \(Y_p=0\) the real part in the leading term is \(O(q^p)\) and the magnitude of its imaginary part is at least half the lower amplitude bound for all sufficiently large \(p\). Multiplication by \(c-\ii s\) in the upper row then bounds \(|Y_n|\) away from zero. The finitely many shorter rows in \eqref{general-insertion} are uniformly bounded. Thus \(|Z|=O_{d,X}(1/n)<|Y_n|\) for all sufficiently large \(n\).

\subsection{The middle interval}
Choose \(X=X(d)\ge16\) so that \(d e^{-x}\rho^{d+4}\le1/32\) for every \(x\ge X\); this is possible by exponential decay. On \(X\le x\le3\sqrt p\), the telescoping pole bound from \eqref{bias} gives \(|V_j|\le\rho^2/2\) for all \(j\), including \(j=0\). Coefficientwise absolute values of \(B\) are bounded by the coefficients of \(\exp(\sum |b_a|z^a)\). Convolving with \(L\) yields
\[
 |Y_j-V_j|\le(\rho^2/2)(e^v-1),\qquad
 v=\sum |b_a|\rho^a\le dr\rho^d.
\]
Since \(v<1/2\), this gives \(\rho |Y_j-V_j|\le dr\rho^{d+3}\le1/32\), and \(|Y_j|\le\rho^2\). At a zero of \(Y_p\) we have \(|V_p|<s/32\). The phase gap \eqref{center} applies because \(ns^2\ge1\). Consequently \(|Y_n|>15s/16\), while \(\rho|Z|\le dr\rho^{d+4}/n\le1/(32n)\) by \eqref{general-insertion}.

\subsection{The far exterior}
Assume \(x\ge3\sqrt p\). For fixed \(d\) and all sufficiently large \(n\) this implies \(x\ge(4d+12)\log n\). Let \(W=\max_{h\ge1}h^p r^h\) and \(S_j=\sum_{h\ge1}(-r)^hh^j\), as in the two-cycle estimate. For \(j=n,p,p-1,\ldots,p-d\), set \(T_j(h)=j![z^j]B(z)e^{hz}\). Geometric expansion of \(L\) near zero gives \(H_j/(1-t)=\sum_{h\ge0}(-r)^h T_j(h)\). For \(h\ge1\) the absolute majorant gives
\[
 |T_j(h)/h^j-1|\le \exp\!\left(\sum_a |b_a|(j/h)^a\right)-1
                 \le e^{dr n^d}-1\le n^{-4}
\]
eventually. To bound the zero label, use the radius \(R_0=e^{x/d}\). Each \(|b_a|R_0^a\le1\), so \(|T_j(0)|\le j!e^d r^{j/d}\). For \(j\ge2d\) we have \(j/d-1\ge j/(2d)\); using \(x/(2d)\ge2\log n\), \(j!\le n^j\), and \(j\ge n-k\), we obtain \(|T_j(0)|\le rn^{-4}\) for sufficiently large \(n\). This additional factor \(r\) allows \(x\) to be arbitrarily large.

The mode and tail estimates in the two-cycle exterior bound concern the unperturbed weights and hold throughout this exterior once \(n\ge512\). Write \(A_j=H_j/(1-t)\). As \(\sum h^p r^h<4W\), \(\sum h^{p+1}r^h<6nW\), and \(r\le W\), the preceding estimates give \(|A_j-S_j|\le5n^{-4}W\) for \(p-d\le j\le p\) and \(|A_n-S_n|\le7n^{-3}W\). At \(A_p=0\), recentering at the stronger neighbor \(h_0\le n\) gives \(|S_n-h_0S_p|>4W/5\), exactly as above. Hence \(|A_n|>3W/4\) eventually, and \(|A_{p-a}|\le5W\) for \(1\le a\le d\). By \eqref{general-insertion},
\[
 |H_n+H'_p|/(1-t)\le5d\,n^d rW<W/4
\]
eventually. This proves the strict sign at every predecessor zero in the exterior.

\subsection{The parameter domain and the original rows}
Choose \(N(k)>k\) above the three eventual bounds. At every size \(j\ge k\), a term with just one excedance must come from a single cycle. Its cycle length is \(j\), its weight is one, and its first Eulerian coefficient is one. Inversion again gives reciprocity. Thus throughout the cube both quotients are monic of the expected degrees, with constant term one and no nonnegative roots. For a prescribed weight vector connect it by a straight segment to the all-one vector of short-cycle weights. The sign test is uniform over this whole segment; the derangement endpoint and the transport argument in Section~\ref{transport} give strict interlacing. Adding back the two simple zeros proves weak interlacing. Setting every short-cycle weight to zero proves the assertion for the minimum cycle length \(k\).

\section{Finite rows at zero mark}\label{finite}
We describe the integer inequalities completing the proof for \(4\le n\le511\). Write \(P_n=P_n(t,0)\), \(y=t/(1+t)^2\), and \(P_n=(1+t)^{n-2}G_n(y)\). Put \(G_1=G_2=0,\ G_3=1\). The original insertion recurrence \cite[Theorem~30]{rr} specializes to
\[
 F_n=(n-1)tF_{n-1}+t(1-t)F'_{n-1}+
     \binom{n-1}{2}t(1+t)F_{n-3}\quad(u=0)
\]
with \(F_0=1,\ F_1=F_2=0\). Substitution after dividing by \(t\) gives for \(n\ge4\), with out-of-range coefficients set to zero,
\[
 [y^i]G_n=(i+1)[y^i]G_{n-1}
       +2(n-1-2i)[y^{i-1}]G_{n-1}
       +\binom{n-1}{2}[y^{i-1}]G_{n-3}.
\]
This follows by applying \((1+2(n-3)y)+y(1-4y)\,d/dy\) to the preceding gamma polynomial. The term of largest possible degree is positive: for even \(n\) its shift coefficient is \(2\), and for odd \(n\) the previous row already has that degree. In particular \(\deg G_n=\lfloor(n-2)/2\rfloor\) and \(G_n(0)=1\).

The supplementary certificate records disjoint negative dyadic intervals for every \(G_n\), in increasing order. To verify a row, evaluate its integer polynomial at both ends of each interval using homogeneous Horner evaluation, check opposite signs, and compare successive endpoints by integer arithmetic. There are as many intervals as the degree, so they exhaust the roots and each root is simple. For each even upper index, the upper row starts and ends the alternation of the two rows' intervals. For each odd upper index, the preceding row starts it and the upper row ends it. The recurrence, degree, bracket signs and ordering can be checked without repeating the search for intervals. The 508 pairs needed here use 64,770 brackets and 129,540 endpoint signs. The supplementary file also contains the longer independent check through size 1023. A second implementation formed the ordinary \(t\)-coefficients and inverted the triangular change of basis before evaluating the same brackets.

On \((-1,0)\), \(t\mapsto y\) is a strictly increasing bijection onto \((-\infty,0)\). Each gamma root gives one root on this interval and one reciprocal root on \((-\infty,-1)\). The term of highest gamma degree shows that at \(-1\) the even \(n\) row is nonzero and the odd \(n\) row has a simple root. For an even upper index the upper row lies closest to \(-1\) on either side and the predecessor has the central root; the roles reverse for odd upper indices. Thus the bracket alternation gives strict interlacing of the original quotients, including at the center. Restoring the common simple root at zero completes the proof of Theorem~\ref{main}.

\paragraph{Formal verification.}
\begingroup\raggedright
The repository's \texttt{verification/README.md} maps Theorems~\ref{general}--\ref{four} to the Lean proofs. Under \texttt{LongCycleBounds/}, the entry modules for Theorems~\ref{general}, \ref{main}, and \ref{family} are \texttt{FixedCutoffCubeVerified.lean}, \texttt{FixedThreeAllSizes.lean}, and \texttt{MarkedHalfVerified.lean}, respectively. The \(k=12\) declarations of type \texttt{Prop} for rows \(23,24\) in \texttt{LongCycles/Statement.lean} are auxiliary specifications.
\par\endgroup

\section{Minimum cycle length four}\label{sec:four}
We give the quantitative estimates for Theorem~\ref{four}. Use the notation of Section~\ref{sec:general} with \(d=3\). The two marks can vary throughout the closed square. The three coefficients of \(\log B(z)\) are
\[
 b_1=\frac r{1+r},\qquad
 b_2=\frac{(1-w_2)r}{2(1+r)^2},\qquad
 b_3=\frac{(1-w_3)r(1-r)}{6(1+r)^3}.
\]
They lie in \([0,1/2],[0,1/8],[0,1/24]\) and are also bounded by \(r,r/2,r/6\). In \eqref{general-insertion}, \(|C_1|,|C_2|,|C_3|\) satisfy the same latter bounds, since \(|1-4r+r^2|\le(1+r)^4\) on \([0,1]\). In particular the sign of \(w_3-w_2\) is unrestricted. In the three estimates below \(n\ge512\), \(p=n-1\), and the smallest coefficient index used is \(n-4\ge508\).

\paragraph{Compact interval.}
Write \(D=|B(\zeta)|\). Since \(\Re(\zeta^3)=x^3-3\pi^2x\ge-2\pi^3\), one has \(D\ge\exp(-\pi^2/8-\pi^3/12)>e^{-5}\).
The contour argument of \eqref{poles} bounds the absolute error in \(Y_j\) after retaining the nearest pole pair by \(170\exp(\sum b_i R^i)(\rho/R)^j\). The following upper bounds hold in the indicated order:
\[
\begin{array}{c|ccc}
x & R & \rho/R & \sum b_iR^i\\ \hline
{[0,4]}&8&7/10&34\\
{[4,8]}&11&9/10&6\\
{[8,16]}&18&24/25&1
\end{array}
\]
For the radii and ratios square both sides and use \(3<\pi<16/5\). For the first cost use the absolute bounds on \(b_i\); on the other intervals use \(r(R+R^2/2+R^3/6)\) with \(r\le e^{-4},e^{-8}\), respectively. At \(j\ge508\) each error is below \(10^{-6}<sD/16\), where \(s=\pi/\rho>3/17\). At \(Y_p=0\) the real part of \(B(\zeta)(\rho/\zeta)^n\) has absolute value below \(sD/32\). Rotating by \(\rho/\zeta\) thus gives \(|Y_n|>3sD/2\), while all three shorter rows have absolute value below \(3D\). For \(2\le h\le4\) one has \(r\rho^h\le\pi^h\), because the logarithmic derivative of \(e^{-x}\rho^h\) is at most \(-1+h/(2\pi)<0\). Hence
\[
 |Z|\le\frac{3D}{512}(\pi^2+\pi^3/2+\pi^4/6)
 < \frac9{34}D < \frac32sD .
\]

\paragraph{Middle interval.}
For \(16\le x\le3\sqrt p\), \(\rho\le33x/32\) and \(r\rho^6\le(33/2)^6e^{-16}<16/5\). Indeed \(x^6 e^{-x}\) decreases here. The tilt cost \(v=\sum b_i\rho^i\) and the drift sum \(J=\sum |C_i|\rho^{i+1}\) satisfy \(v\le(21/100)r\rho^3,\ J\le(21/100)r\rho^4\), using \(\rho\ge16\). The logistic coefficients satisfy \(|V_j|\le\rho^2/2\) for every \(j\), by the pole mass estimate (including \(V_0\)). Convolution gives \(|Y_j-V_j|\le(\rho^2/2)(e^v-1)\le\rho^2v<s/4\) as \(v\rho^3<27/40<\pi/4\). At \(Y_p=0\), \eqref{center} applies since \(ns^2\ge1\) and \(|V_p|<s/4\le1/4\); it yields \(|V_n-(x/\rho)V_p|>s\), so \(|Y_n|>s/2\). Every shorter row is bounded by \(3\rho^2/5\). Consequently \(n|Z|\le(3/5)(21/100)r\rho^6<1/2\le ns/2\).

\paragraph{Exterior.}
For \(x\ge3\sqrt p\) we have \(r n^8\le1\) as in the two-cycle case. In the expansion \(A_j=H_j/(1-t)=\sum_{h\ge0}(-r)^h T_j(h)\), one has \(T_j(h)=j![z^j]B(z)e^{hz}\). For \(h\ge1,\ j\le n\), the relative error \(|T_j(h)/h^j-1|\) is bounded by \(\exp(r(n+n^2/2+n^3/6))-1<n^{-4}\), using falling factorials bounded by the corresponding powers of \(j\). For the zero label take \(R_0=e^{x/3}\ge n^2\); Cauchy's estimate gives \(|T_j(0)|\le j!e^3 R_0^{-j}\le r n^{-4}\), because \(j+5\le2(j-3)\) for the indices here. Write \(W=\max_{h\ge1}h^p r^h\), \(S_j=\sum_{h\ge1}(-r)^h h^j\).
For \(p-3\le j\le n\) the moment estimates from the preceding exterior argument give \(|A_j-S_j|\le W/(8n)\): the positive-label sums are bounded by \(6nW\) for \(j\le p+1\), and \(r\le W\). At \(A_p=0\), subtracting the stronger neighbor index times \(S_p\) as before gives \(|A_n|>W/4\). For the three preceding indices \(|S_j|\le W\) by the alternating unimodal sum, hence \(|A_j|<2W\). The insertion in raw coordinates is bounded by \(6W r n^3<W/4\). This completes the derivative sign test on the square. The same path to the all-one derangement row and the same reflection argument now give strict interlacing of the quotients throughout that square.

\paragraph{Remaining original rows.}
At \(w_2=w_3=0\) the original row recurrence has short-cycle term \(\binom{n-1}{3}t(1+4t+t^2)E_{n-4,4}\), since the six cyclic orderings on four given letters have excedance polynomial \(t+4t^2+t^3\). In the gamma coordinates of Section~\ref{finite}, start with \(G_1=G_2=G_3=0,\ G_4=1+2y\). For \(n\ge5\),
\[
 G_n=(1+2(n-3)y)G_{n-1}+y(1-4y)G'_{n-1}
       +\binom{n-1}{3}y(1+2y)G_{n-4}.
\]
The constant term is one and the degree is \(\lfloor(n-2)/2\rfloor\) by the same top-coefficient argument. The supplementary dyadic certificate checks the 507 adjacent pairs between rows 4 and 511; counting the base row, these use 64,770 disjoint within-row brackets and 129,540 integer endpoint signs. One checker uses the gamma recurrence and the other reconstructs the ordinary rows before extracting their gamma coefficients. Degrees, interval order, reflection and the central root are handled exactly as in Section~\ref{finite}. This proves Theorem~\ref{four}.

\begingroup\raggedright
Theorem~\ref{four} uses \texttt{LongCycleBounds/FixedFourAllSizes.lean} for the original rows and \texttt{LongCycleBounds/FixedFourVerified.lean} for the weighted result. The finite witnesses and their Python checkers are in \texttt{certificates/k3/} and \texttt{certificates/k4/}. The local Lean build, kernel replay, and axiom audit records are indexed in \texttt{verification/README.md}.
\par\endgroup

\paragraph{Acknowledgment.}
ChatGPT assisted with the analytic derivations, the implementation and independent replay of the formal proofs, and preparation of the manuscript.
\begin{thebibliography}{9}
\bibitem{alex}
P.~Alexandersson, \emph{Real-rooted polynomials: open problems},
\href{https://www.symmetricfunctions.com/realRootedOpenProblems.htm\#permutations-with-long-cycles}{Permutations with long cycles},
The Symmetric Functions Catalog, 2026.
\bibitem{programcomputer}
ProgramComputer, \emph{Consecutive excedance zeros of permutations with a minimum cycle length}, November 2026,
\href{https://github.com/ProgramComputer/open-math-research/blob/4da48b72ed8f92f96c2faf60e2d63d253b3b507a/topics/long-cycle-excedances/paper.tex}{Theorem~1 and the uniform bound}.
\bibitem{brenti}
F.~Brenti, \emph{Permutation enumeration, symmetric functions, and
unimodality}, Pacific J. Math. \textbf{157} (1993), 1--28,
\href{https://msp.org/pjm/1993/157-1/pjm-v157-n1-p01-s.pdf}{Theorems~3.2--3.4 and Problem~5.5}.
\bibitem{rr}
I.-N.~M.~Rakotoson and F.~Rakotondrajao,
\emph{Associated Stirling Eulerian numbers and their \(q\)-analogues},
J. Integer Seq. \textbf{28} (2025),
\href{https://cs.uwaterloo.ca/journals/JIS/VOL28/Rakotoson/rakoto12.html}{Article 25.5.7},
Theorem~30.
\bibitem{dlmf}
NIST Digital Library of Mathematical Functions,
\href{https://dlmf.nist.gov/4.36.E2}{Eq.~4.36.2}.
\end{thebibliography}
\end{document}
